\documentclass[handout,aspectratio=169,10pt]{beamer}
%\documentclass[aspectratio=169,10pt]{beamer}  % Alternative for animations

\input{misc/style}

\title{Title}
\date{\today}
\author{Presenter Name}
\institute[]{Institution Name}
% Add an authorized organization logo if needed.

\begin{document}
\maketitle

%%%%%%%%%%%%%%%%%%%%%%%%%%%%%%%%%%%%%%%%
% Roadmap

\begin{frame}{Roadmap}
\tableofcontents
\end{frame}

%%%%%%%%%%%%%%%%%%%%%%%%%%%%%%%%%%%%%%%%
% Main content
%
% Visual Primacy Principle: every content frame should lead with a
% figure, table, or equation. Text is a caption, not the message.
% Two-column layouts give the visual ≥0.6\textwidth; the text column
% carries 2-4 short labels, not prose.

\section{First Section}

\begin{frame}{Descriptive Frame Title}
    \begin{columns}[T]
        \begin{column}{0.6\textwidth}
            \begin{block}{Cash balance identity}
            \[C_{\mathrm{end}} = C_{\mathrm{start}} + R - P\]
            \end{block}
        \end{column}
        \begin{column}{0.36\textwidth}
            \footnotesize
            \begin{itemize}
                \item $R$: receipts; $P$: payments
                \item Check the closing balance
            \end{itemize}
        \end{column}
    \end{columns}
\end{frame}

%%%%%%%%%%%%%%%%%%%%%%%%%%%%%%%%%%%%%%%%
% References

\appendix

\section*{References}
\begin{frame}[allowframebreaks]{References}
    \bibliographystyle{plainnat}
    \bibliography{misc/ref}
\end{frame}

\end{document}
